% Options for packages loaded elsewhere
\PassOptionsToPackage{unicode}{hyperref}
\PassOptionsToPackage{hyphens}{url}
\PassOptionsToPackage{dvipsnames,svgnames,x11names}{xcolor}
%
\documentclass[
  11pt,
]{article}
\usepackage{amsmath,amssymb}
\usepackage{iftex}
\ifPDFTeX
  \usepackage[T1]{fontenc}
  \usepackage[utf8]{inputenc}
  \usepackage{textcomp} % provide euro and other symbols
\else % if luatex or xetex
  \usepackage{unicode-math} % this also loads fontspec
  \defaultfontfeatures{Scale=MatchLowercase}
  \defaultfontfeatures[\rmfamily]{Ligatures=TeX,Scale=1}
\fi
\usepackage{lmodern}
\ifPDFTeX\else
  % xetex/luatex font selection
\fi
% Use upquote if available, for straight quotes in verbatim environments
\IfFileExists{upquote.sty}{\usepackage{upquote}}{}
\IfFileExists{microtype.sty}{% use microtype if available
  \usepackage[]{microtype}
  \UseMicrotypeSet[protrusion]{basicmath} % disable protrusion for tt fonts
}{}
\makeatletter
\@ifundefined{KOMAClassName}{% if non-KOMA class
  \IfFileExists{parskip.sty}{%
    \usepackage{parskip}
  }{% else
    \setlength{\parindent}{0pt}
    \setlength{\parskip}{6pt plus 2pt minus 1pt}}
}{% if KOMA class
  \KOMAoptions{parskip=half}}
\makeatother
\usepackage{xcolor}
\usepackage[margin=1in]{geometry}
\usepackage{longtable,booktabs,array}
\usepackage{calc} % for calculating minipage widths
% Correct order of tables after \paragraph or \subparagraph
\usepackage{etoolbox}
\makeatletter
\patchcmd\longtable{\par}{\if@noskipsec\mbox{}\fi\par}{}{}
\makeatother
% Allow footnotes in longtable head/foot
\IfFileExists{footnotehyper.sty}{\usepackage{footnotehyper}}{\usepackage{footnote}}
\makesavenoteenv{longtable}
\setlength{\emergencystretch}{3em} % prevent overfull lines
\providecommand{\tightlist}{%
  \setlength{\itemsep}{0pt}\setlength{\parskip}{0pt}}
\setcounter{secnumdepth}{-\maxdimen} % remove section numbering
\DeclareUnicodeCharacter{220E}{\ensuremath{\square}}
\ifLuaTeX
  \usepackage{selnolig}  % disable illegal ligatures
\fi
\IfFileExists{bookmark.sty}{\usepackage{bookmark}}{\usepackage{hyperref}}
\IfFileExists{xurl.sty}{\usepackage{xurl}}{} % add URL line breaks if available
\urlstyle{same}
\hypersetup{
  pdfauthor={Author information to be supplied before submission},
  colorlinks=true,
  linkcolor={Maroon},
  filecolor={Maroon},
  citecolor={Blue},
  urlcolor={Blue},
  pdfcreator={LaTeX via pandoc}}

\title{A Width-One Generalized-Davenport Corridor in \(C_5^3\) and a Rank-Forcing Obstruction Phase}
\author{Author information to be supplied before submission}
\date{Review source - 25 August 2026}

\begin{document}
\maketitle

\hypertarget{abstract}{%
\section{Abstract}\label{abstract}}

Let \(D_k(G)\) be the least length that forces \(k\) pairwise disjoint
nonempty zero-sum subsequences. We show that the generalized Davenport
constants of \(C_5^3\) lie in the width-one corridor

\[
5k+10\le D_k(C_5^3)\le5k+11
\qquad(k\ge4).
\]

If \(D_4(C_5^3)=30\), then the lower line is exact for every \(k\ge2\).
To analyze the remaining alternative, we study a hypothetical total-zero
sequence of length \(31\) with no nonempty zero-sum subsequence of
length at most five. For saturated short-free sequences over elementary
abelian groups of odd prime exponent, we prove a defect certificate that
excludes multiplicity \(p-2\). In the case \(p=5\), only multiplicities
\(1,2,4\) remain. If \(s\) is support size and \(c_i\) counts
multiplicity \(i\), then

\[
c_2=31-s-3c_4,
\qquad
c_1=2s-31+2c_4.
\]

The subsequence of points of multiplicity at least two has length
\(62-2s-2c_4\). Using \(\eta(C_5^2)=13\), we prove that this stratum
must span rank three whenever \(s+c_4\le24\). We also prove an exact
identity coupling support deficit, repetitions inside zero-sum atoms,
and overlaps between atom supports. Together with a theorem-grade
exclusion through support ten, these results isolate the residual
obstruction phase without deciding it.

The exact value remains \(D_4(C_5^3)\in\{30,31\}\). A larger internal
search through support 22 is reported only as bounded computational
evidence and is not used as theorem authority.

\textbf{Keywords:} generalized Davenport constants; zero-sum sequences;
elementary abelian groups; short zero sums; inverse zero-sum problems

\hypertarget{introduction}{%
\section{1. Introduction}\label{introduction}}

Generalized Davenport constants measure the sequence length required to
force several pairwise disjoint zero sums. Rank-two groups exhibit
strong eventual structure, whereas rank-three groups can depart from the
expected arithmetic progression.

For \(C_5^3\), the exact values \(D_2=20\) and \(D_3=25\), together with
known short-zero-sum thresholds, reduce the next value to

\[
D_4(C_5^3)\in\{30,31\}.
\]

This one-bit ambiguity controls the tail: the lower candidate \(D_4=30\)
forces \(D_k=5k+10\) for every \(k\ge2\). The upper candidate would
yield a rigid length-31 total-zero sequence with no zero sum of length
at most five.

We combine the recurrence with symbolic analysis of that hypothetical
obstruction. The purpose is twofold: to obtain a rigorous width-one
theorem for the full tail and to identify exactly where existing rank
reductions cease to be automatic. Computation is kept in a separate
evidence class.

Our contributions are:

\begin{enumerate}
\def\labelenumi{\arabic{enumi}.}
\tightlist
\item
  the width-one corridor for every \(k\ge4\);
\item
  the conditional exact tail if \(D_4=30\);
\item
  an odd-prime saturation-defect lemma excluding multiplicity \(p-2\);
\item
  a theorem-grade support-at-least-eleven result for the length-31
  obstruction;
\item
  an exact multiplicity parameterization for every remaining support
  stratum;
\item
  the rank-forcing phase \(s+c_4\le24\) for the high-multiplicity
  stratum; and
\item
  an identity coupling support deficit to internal atom repetition and
  cross-atom overlap.
\end{enumerate}

\hypertarget{definitions-and-known-inputs}{%
\section{2. Definitions and known
inputs}\label{definitions-and-known-inputs}}

For a finite abelian group \(G\), \(D_k(G)\) is the least integer \(n\)
such that every sequence of length at least \(n\) over \(G\) contains
\(k\) pairwise disjoint nonempty zero-sum subsequences. Let
\(s_{\le\ell}(G)\) denote the least length forcing a nonempty zero-sum
subsequence of length at most \(\ell\).

We use the generalized-Davenport recurrence framework {[}1{]}

\[
D_{k+1}(G)
\le\max\{D_k(G)+\ell, \ s_{\le\ell}(G)-1\},
\]

and the known lower bound \(D_k(C_5^3)\ge5k+10\) in the range considered
here. The exact registered inputs are

\[
D_2=20, \qquad D_3=25, \qquad s_{\le6}=24,
\]

while \(s_{\le5}(C_5^3)=\eta(C_5^3)=33\). We also use the rank-two value
\(\eta(C_5^2)=13\).

\hypertarget{a-width-one-tail}{%
\section{3. A width-one tail}\label{a-width-one-tail}}

Taking \(\ell=6\) gives

\[
D_4\le\max\{25+6,24-1\}=31.
\]

Together with the lower bound, this proves \(D_4\in\{30,31\}\).

\textbf{Theorem 1.} For every \(k\ge4\),

\[
5k+10\le D_k(C_5^3)\le5k+11.
\]

\textbf{Proof.} The base case is \(D_4\le31\). Apply the recurrence with
\(\ell=5\) and \(s_{\le5}-1=32\). If \(D_k\le5k+11\), then both terms in
the maximum are at most \(5(k+1)+11\). The known lower bound supplies
the lower line. ∎

\textbf{Theorem 2.} If \(D_4(C_5^3)=30\), then

\[
D_k(C_5^3)=5k+10
\qquad(k\ge2).
\]

\textbf{Proof.} The statement holds for \(k=2,3,4\). Applying the same
recurrence with \(\ell=5\) propagates the upper bound \(5(k+1)+10\),
which matches the known lower bound. ∎

There is no converse here: the assumption \(D_4=31\) does not by itself
prove that later terms lie on the upper line.

\hypertarget{the-hypothetical-length-31-obstruction}{%
\section{4. The hypothetical length-31
obstruction}\label{the-hypothetical-length-31-obstruction}}

The upper candidate yields a sequence \(S\) over \(C_5^3\) satisfying

\[
|S|=31, \qquad \sigma(S)=0,
\]

with no nonempty zero-sum subsequence of length at most five. Every
element has multiplicity at most four. In the support case relevant
here, the established extension dichotomy forces saturation: appending
any support element destroys short-freeness.

\textbf{Theorem 3 (saturation defect).} Let \(p\) be odd and let \(S\)
be a saturated \(p\)-short-free sequence over an elementary abelian
exponent-\(p\) group. If a nonzero element \(x\) has multiplicity
\(m<p\), then there is a subsequence \(R\) such that

\[
|R|\le p-1-m,
\qquad x\notin\operatorname{supp}(R),
\qquad \sigma(R)=-(m+1)x.
\]

\textbf{Proof.} Append one further copy of \(x\). Saturation supplies a
zero-sum subsequence of length at most \(p\) that uses the appended
occurrence. It must also use every original copy of \(x\); otherwise
replacing the appended copy by an unused original occurrence would give
a short zero sum already contained in \(S\). Removing the \(m+1\) copies
of \(x\) leaves the required \(R\). ∎

\textbf{Corollary 4.} Multiplicity \(p-2\) is impossible.

Indeed, the defect subsequence would have length at most one and sum
\(x\), forcing a forbidden additional copy of \(x\). For \(p=5\),
multiplicity three is absent. Singleton and double points also acquire
defect certificates of length at most three and two, respectively.

\hypertarget{exact-multiplicity-grammar}{%
\section{5. Exact multiplicity
grammar}\label{exact-multiplicity-grammar}}

Let \(s=|\operatorname{supp}(S)|\), and let \(c_1, c_2, c_4\) count
points of multiplicity one, two, and four. Then

\[
c_1+c_2+c_4=s,
\qquad
c_1+2c_2+4c_4=31.
\]

Writing \(c_4=j\) gives

\[
c_2=31-s-3j,
\qquad
c_1=2s-31+2j.
\]

Thus \(c_1\) is odd, and nonnegativity determines the complete
admissible range for every support size. For example, the support-23
patterns are

\[
1^{15}2^8, \qquad 1^{17}2^5 4, \qquad 1^{19}2^2 4^2,
\]

and the support-24 patterns are

\[
1^{17}2^7, \qquad 1^{19}2^4 4, \qquad 1^{21}2 4^2.
\]

\hypertarget{the-rank-forcing-phase}{%
\section{6. The rank-forcing phase}\label{the-rank-forcing-phase}}

Let \(H\) be the subsequence consisting of all points of multiplicity
two or four. Its length is

\[
|H|=2c_2+4c_4=62-2s-2c_4.
\]

If \(H\) were contained in a subgroup of rank at most two, it would be a
5-short-free sequence over \(C_5^2\). The value \(\eta(C_5^2)=13\)
forces a short zero sum once \(|H|\ge13\). Since \(|H|\) is even, the
contradiction begins at \(|H|\ge14\).

\textbf{Theorem 5 (rank-forcing phase).} The high-multiplicity stratum
spans rank three whenever

\[
s+c_4\le24.
\]

\textbf{Proof.} The condition \(|H|\ge14\) is equivalent to
\(62-2s-2c_4\ge14\), hence to \(s+c_4\le24\). ∎

The theorem gives the following exact boundary:

\begin{longtable}[]{@{}rrl@{}}
\toprule\noalign{}
Support & \(c_4\) & Consequence from Theorem 5 \\
\midrule\noalign{}
\endhead
\bottomrule\noalign{}
\endlastfoot
23 & 0 or 1 & rank three forced \\
23 & 2 & threshold is silent \\
24 & 0 & rank three forced \\
24 & 1 or 2 & threshold is silent \\
at least 25 & any admissible value & threshold is silent \\
\end{longtable}

Silence is not a converse. Outside the phase, rank three may follow from
a different argument.

\hypertarget{theorem-grade-low-support-exclusion-and-bounded-evidence}{%
\section{7. Theorem-grade low-support exclusion and bounded
evidence}\label{theorem-grade-low-support-exclusion-and-bounded-evidence}}

Signed symbolic reductions combined with two exact state representations
exclude every support size through ten.

\textbf{Theorem 6.} Every length-31 total-zero 5-short-free sequence
over \(C_5^3\), if one exists, has support at least eleven.

A larger internal computation found no such sequence through support 22.
That search has not received an external independent replay and is not
used in any proof above or below. Its bounded conclusion is useful for
prioritizing the next case split, but it does not establish support at
least 23 as a theorem.

\hypertarget{atom-repetition-and-overlap}{%
\section{8. Atom repetition and
overlap}\label{atom-repetition-and-overlap}}

Under the conditional assumption \(D_4=31\), extremal factorization
arguments force one of the atom-length types

\[
(6,6,6,13)
\qquad\text{or}\qquad
(6,6,7,12).
\]

Let \(S=U_1\cdots U_r\) be an occurrence-disjoint atom factorization.
For each group element \(g\), let \(r_g\) be the number of atom supports
containing \(g\), and define the internal deficit

\[
\delta_i=|U_i|-|\operatorname{supp}(U_i)|.
\]

\textbf{Theorem 7 (repetition-overlap identity).} Every
occurrence-disjoint factorization satisfies

\[
|S|-|\operatorname{supp}(S)|
=\sum_i\delta_i
+\sum_{g\in\operatorname{supp}(S)}(r_g-1).
\]

\textbf{Proof.} Summing the atom support sizes counts every global
support point \(r_g\) times. Hence

\[
\sum_i\delta_i=|S|-\sum_g r_g,
\]

while cross-atom overlap equals

\[
\sum_g r_g-|\operatorname{supp}(S)|.
\]

Adding the two identities proves the claim. ∎

Thus any independently established lower support bound \(s_0\) gives a
total internal-repetition plus cross-overlap budget of at most
\(31-s_0\). The proven value \(s_0=11\) gives budget 20. If the bounded
support-22 computation is later independently certified, the budget
sharpens to eight; that sharper value is conditional and not used as
theorem authority here.

\hypertarget{the-residual-obstruction-phase}{%
\section{9. The residual obstruction
phase}\label{the-residual-obstruction-phase}}

A decisive next analysis should impose simultaneously:

\begin{enumerate}
\def\labelenumi{\arabic{enumi}.}
\tightlist
\item
  the multiplicity pattern from Section 5;
\item
  the rank status from Theorem 5;
\item
  the saturation-defect certificate at every singleton and double point;
\item
  one of the two conditional atom-length types;
\item
  the repetition-overlap identity; and
\item
  total sum zero with exclusion of zero sums of lengths one through
  five.
\end{enumerate}

This intersection is substantially smaller than an unstructured support
search. A surviving sequence would require independent factorization
verification before establishing \(D_4=31\). A complete, independently
replayed infeasibility proof for the length-31 target would establish
the lower candidate \(D_4=30\).

\hypertarget{relation-to-prior-work}{%
\section{10. Relation to prior work}\label{relation-to-prior-work}}

The generalized-Davenport recurrence, lower-bound framework, extremal
factorization language, and eventual-linearity context are established
results {[}1{]}. The \(C_0\) framework and short-zero-sum localization
are also prior work {[}2{]}, as are the inverse zero-sum program {[}3{]}
and rank-two inverse results {[}4{]}.

The residual contribution is the width-one synthesis for \(C_5^3\), the
saturation-defect specialization, the exact multiplicity grammar, the
rank-forcing phase, and the coupling of support deficit with atom
repetition and overlap. These are structural compression results; they
do not decide the final extremal bit.

\hypertarget{reproducibility-and-limitations}{%
\section{11. Reproducibility and
limitations}\label{reproducibility-and-limitations}}

The corridor, defect lemma, and support-ten exclusion are separately
bound to exact records. An independent verifier enumerates all
multiplicity patterns for supports 8-31, checks the rank-phase
equivalence, reproduces the support-23 and support-24 tables, and
validates the atom identity on independently generated factorizations.
The displayed proofs carry all-parameter authority.

The exact value of \(D_4(C_5^3)\) remains open, as does the associated
length-31 extremal-spectrum statement. The support-22 frontier is
bounded computation pending external replay. The rank phase is one-way,
and the atom identity is a compression principle rather than a
contradiction. No additional internal implementation can substitute for
an independently auditable final obstruction proof.

\hypertarget{conclusion}{%
\section{12. Conclusion}\label{conclusion}}

The generalized Davenport constants of \(C_5^3\) remain within one unit
of the line \(5k+10\) for every \(k\ge4\), and the lower choice at
\(k=4\) would fix the entire tail. The hypothetical upper-line
obstruction has substantially more structure than this corridor alone
reveals: its multiplicities have an exact grammar, its high-multiplicity
stratum enters a sharp rank-forcing phase, and its atom factorization
obeys a global repetition-overlap budget.

These results explain where the residual difficulty begins and define a
smaller intersection problem for the unresolved extremal bit. They do
not claim to have resolved it.

\hypertarget{data-and-code-availability}{%
\section{Data and code availability}\label{data-and-code-availability}}

The verification package contains exact multiplicity enumeration,
rank-phase checks, atom-identity tests, and separately labeled
bounded-search records. A permanent archival identifier should be added
before final submission.

\hypertarget{references}{%
\section{References}\label{references}}

\begin{enumerate}
\def\labelenumi{\arabic{enumi}.}
\tightlist
\item
  M. Freeze and W. A. Schmid, ``Remarks on a Generalization of the
  Davenport Constant,'' \emph{Discrete Mathematics} \textbf{310},
  3373-3389 (2010). DOI: 10.1016/j.disc.2010.07.028
\item
  Y. Fan, W. Gao, G. Wang, Q. Zhong, and J. Zhuang, ``On Short Zero-Sum
  Subsequences of Zero-Sum Sequences,'' \emph{Electronic Journal of
  Combinatorics} \textbf{19}(3), P31 (2012). DOI: 10.37236/2602
\item
  W. Gao, A. Geroldinger, and W. A. Schmid, ``Inverse Zero-Sum
  Problems,'' \emph{Acta Arithmetica} \textbf{128}, 245-279 (2007). DOI:
  10.4064/aa128-3-5
\item
  Q. Zhong, ``On the Inverse Problem of the \(k\)-th Davenport Constants
  for Groups of Rank 2,'' \emph{Combinatorica} \textbf{45}, article 31
  (2025). DOI: 10.1007/s00493-025-00153-3
\end{enumerate}

\end{document}
